\documentclass[11pt]{article}
\usepackage{amsmath,amsthm,amssymb}
\usepackage[margin=1in]{geometry}
\usepackage{booktabs}
\usepackage{array}
\usepackage{stmaryrd}
\usepackage[colorlinks=true,linkcolor=blue,citecolor=blue,urlcolor=blue]{hyperref}

\newtheorem{theorem}{Theorem}
\newtheorem{proposition}[theorem]{Proposition}
\newtheorem{corollary}[theorem]{Corollary}
\newtheorem{lemma}[theorem]{Lemma}
\theoremstyle{definition}
\newtheorem{definition}[theorem]{Definition}
\newtheorem{remark}[theorem]{Remark}
\newtheorem{example}[theorem]{Example}

\newcommand{\Poly}{\mathsf{Poly}}
\newcommand{\PolyZ}{\mathsf{Poly}_{\mathbb{Z}}}
\newcommand{\SPoly}{\mathsf{SPoly}}
\newcommand{\Cat}{\mathsf{Cat}}
\newcommand{\Set}{\mathsf{Set}}
\newcommand{\Comon}{\mathsf{Comon}}
\newcommand{\Ob}{\mathrm{Ob}}
\newcommand{\Hom}{\mathrm{Hom}}
\newcommand{\N}{\mathbb{N}}
\newcommand{\Z}{\mathbb{Z}}
\newcommand{\id}{\mathrm{id}}
\DeclareMathOperator{\rank}{rank}

\title{Parallelizability does not port to directed containers\\
\large A tangent dividing line, and the out-degree invariant blind to composition}
\author{MacBeth\thanks{Kodamai. Source, this note, and the Lean certification live at
\texttt{github.com/scotmacbeth/work-in-progress}; see the content commit recorded on this page.
Prepared for the Strathclyde mini-course \emph{Tangent Categories and Containers}
(13--16 October 2026) as the cautionary companion to the derivative-uniqueness arc.}\\
\small Kodamai}
\date{6 October 2026\\[2pt]
\small content commit \texttt{e154282}}

\begin{document}
\maketitle

\begin{abstract}
Lanfranchi proves that an object of a Cartesian tangent category is \emph{parallelizable} exactly
when it carries the structure of a Lie group object. Directed containers are small categories, and
the one-hole derivative $\partial$ of Abbott, Altenkirch, Ghani and McBride is a tangent structure
on polynomial functors, so it is tempting to conjecture that a container is parallelizable
precisely when it is a groupoid. We show this port is ill-posed, and isolate why. Lanfranchi's
group objects live in a \emph{Cartesian} tangent category, whereas the derivative's tangent
structure is the Weil base change $T=(-)\otimes W$, which is representable/monoidal and not
Cartesian: $T$ fails to preserve the categorical product by an explicit second-order term, and
with the product as trivial bundle only the constant functors are parallelizable; the genuine
Cartesian home of $\partial$ is a generalized Cartesian differential category, where every object
is parallelizable and a container is a morphism, not an object. What survives is one sharp
invariant: the tangent bundle $Tp$ has, over a point of shape $a$, a fibre of rank equal to the
out-degree $|E_a|$, so all fibres agree iff the category is out-degree homogeneous. This
``tangent-regularity'' depends only on the underlying polynomial functor, hence is provably blind
to composition, and it refutes the groupoid conjecture in both directions. Connected groupoids are
tangent-regular --- the faithful shadow of ``Lie group $\Rightarrow$ parallelizable'' --- but
strictly more categories are. The tangent geometry of containers reads the shape of a category,
not its composition.
\end{abstract}

\section{Introduction}\label{sec:intro}

\paragraph{A derivative, a tangent structure, and a tempting analogy.}
Containers \cite{AAGM2003,GambinoKock2013} have a derivative. The one-hole-context operation
$\partial$ of Abbott, Altenkirch, Ghani and McBride sends a container to the container of its
one-hole contexts and obeys the laws of a derivative. Two companion notes make this sharp: the
category $\SPoly$ of finitary polynomial functors is a Cartesian differential category whose
differential combinator \emph{is} $\partial$ \cite{MacBethTangent}, and the attached first-order
tangent structure --- substitute a dual number and read off the $\varepsilon$-coefficient,
$Tp = p(y+\varepsilon v)\bmod\varepsilon^2 = p + \varepsilon v\cdot\partial p$ --- is essentially
forced \cite{MacBethUniqueness}. Directed containers, moreover, are not just functors: by
Ahman and Uustalu they are \emph{exactly small categories}, the comonoids for the composition
monoidal product $\lhd$ on polynomial functors \cite{AhmanUustalu2016}.

So containers carry two geometries at once --- a tangent structure and a category structure ---
and a recent theorem invites us to relate them. Lanfranchi \cite{Lanfranchi2026} proves, for an
object $G$ of a \emph{Cartesian} tangent category, a Lie-theoretic correspondence whose fourth
equivalent clause is that $G$ is \emph{parallelizable}: its tangent bundle trivializes,
$TG \cong G\times m$ over $G$, for a differential object $m$ (Definition~4.18, Theorem~4.20). In
smooth manifolds parallelizability is a rare and structured condition --- Lie groups are
parallelizable, generic manifolds are not. The analogy writes itself: a directed container is a
small category; which categories are parallelizable, and is the answer ``the groupoids''? A
groupoid is the categorical avatar of a group, and ``Lie group $\Rightarrow$ parallelizable'' is
the classical half of Lanfranchi's theorem.

\paragraph{The caveat, and the result.}
There is a structural mismatch to clear first. Lanfranchi's group objects are internal groups for
the \emph{Cartesian} product; directed containers are comonoids for \emph{composition} $\lhd$.
These are different monoidal structures on $\Poly$, and the derivative's tangent functor is
Cartesian for neither. This note shows the mismatch is fatal to the literal port, and extracts
what remains. We prove two things, stated here and made precise in
Sections~\ref{sec:fibre}--\ref{sec:survivor}.

\begin{theorem}[The Cartesian port is ill-posed]\label{thm:I-intro}
There is no Cartesian tangent category on polynomial functors in which the derivative $\partial$
is the tangent bundle and Lanfranchi's parallelizability is a nontrivial condition.
\begin{itemize}
\item On $(\Poly,\times)$ the tangent functor $T=(-)\otimes W$ does not preserve the categorical
product --- it leaves a nonzero second-order ($\varepsilon^2$) term --- so $(\Poly,\times,T)$ is
not a Cartesian tangent category; and with $\times$ as the trivial bundle the \emph{only}
parallelizable objects are the constants (even the terminal object $y$ fails).
\item The tangent structure of $\partial$ is genuinely the \emph{representable/monoidal} base
change $(-)\otimes W$, not a Cartesian one.
\item The honest Cartesian home of $\partial$ is the arity generalized Cartesian differential
category $\SPoly$, where every object is parallelizable (vacuous) and a container is a
\emph{morphism}, not an object.
\end{itemize}
\end{theorem}

\begin{theorem}[The surviving invariant, and the groupoid refutation]\label{thm:II-intro}
The tangent bundle $Tp$ of a container $p=\sum_{a\in\Ob}y^{E_a}$ has, over a point of shape $a$,
a fibre that is the free module on the direction set $E_a$, of rank $|E_a|$. Hence the fibres are
all isomorphic to one module iff the out-degree $a\mapsto|E_a|$ is constant. This condition ---
``tangent-regularity'' --- depends only on the underlying polynomial functor, hence cannot be
equivalent to ``groupoid'': the monoid $\{1,e\}$ with $e^2=e$ is tangent-regular but is not a
groupoid, and the groupoid $\Z/2\sqcup\Z/3$ is not tangent-regular. Connected groupoids are
tangent-regular, the faithful shadow of Lanfranchi's theorem, but tangent-regularity is strictly
weaker.
\end{theorem}

\paragraph{Method and moral.}
The whole analysis turns on one calculation (\S\ref{sec:fibre}): the fibre of the tangent bundle
over a point of shape $a$ is the free module on $E_a$. From this single fact both theorems follow
by elementary means --- the Cartesian obstruction is a cardinality mismatch, and the surviving
invariant is a rank. The moral is a dividing line we have met before in this program, here in its
cautionary form: results about $\partial$ are results about the \emph{shape} of a container --- its
out-degree sequence --- and are blind to the composition $\delta\colon p\to p\lhd p$ that makes it
a category. ``Groupoid'', ``connected'', ``invertible'' are properties of $\delta$; no invariant
of the shape alone can see them. The Lie-theoretic content of parallelizability, if it survives at
all, lives in a strictly finer, composition-sensitive layer that we can describe but not yet
characterize (\S\ref{sec:conclusion}).

None of this is a defect in \cite{Lanfranchi2026}: the theorem there is about Cartesian tangent
categories, and the port fails precisely because the derivative's tangent structure is not one.
The arithmetic core of Theorem~\ref{thm:II-intro} and its composition-blindness are
machine-checked in Lean~4 \cite{MacBethLean}.

\paragraph{Outline.} Section~\ref{sec:prelim} fixes the container, the derivative, the tangent
structure, and Lanfranchi's definition. Section~\ref{sec:fibre} computes the tangent fibre.
Section~\ref{sec:refutation} proves Theorem~\ref{thm:I-intro}. Section~\ref{sec:survivor} proves
Theorem~\ref{thm:II-intro}. Section~\ref{sec:conclusion} states the moral and the one open
refinement.

\section{Preliminaries}\label{sec:prelim}

We work with polynomial functors in one variable and finite support; everything is first order.

\paragraph{Containers as categories.}
A \emph{container} in one variable is a polynomial functor
\[
   p \;=\; \sum_{a\in\Ob} y^{E_a},
\]
with a set $\Ob$ of \emph{shapes} and, for each shape $a$, a set $E_a$ of \emph{directions}; its
extension sends a set $X$ to $\sum_{a}X^{E_a}$, so a point of $p$ over $X$ is a pair
$(a, f\colon E_a\to X)$ \cite{AAGM2003,GambinoKock2013}. The composition product $\lhd$ has unit
$y$ and realizes substitution of functors. By Ahman and Uustalu \cite{AhmanUustalu2016}, the
comonoids in $(\Poly,\lhd,y)$ are exactly small categories: a comonoid structure
$\delta\colon p\to p\lhd p$ on $p$ is a choice of identities and a composition, $\Ob$ is the
object set, and $E_a$ is the set of morphisms out of $a$ --- so $|E_a|$ is the \emph{out-degree}
of the object $a$. We write $\Comon(\Poly,\lhd)\simeq\Cat$ for this equivalence and
$U\colon\Cat\to\Poly$ for the forgetful functor that discards $\delta$ and keeps only the
underlying functor $\sum_a y^{E_a}$.

\paragraph{The derivative and its tangent structure.}
The one-hole derivative \cite{AAGM2003} is
\[
   \partial p \;=\; \sum_{a}\ \sum_{e\in E_a} y^{E_a\setminus\{e\}}
             \;=\; \sum_a |E_a|\cdot y^{\,|E_a|-1}.
\]
Let $W=\Z[\varepsilon]/(\varepsilon^2)$ be the dual-number Weil algebra. The tangent functor is the
Weil base change $T=(-)\otimes W$ \cite{MacBethCartan}; on the first-order (species) layer it is
computed by substituting the infinitesimal and truncating,
\begin{equation}\label{eq:T}
   Tp \;=\; p(y+\varepsilon v)\bmod\varepsilon^2 \;=\; p + \varepsilon v\cdot\partial p,
\end{equation}
where $v$ is the tangent coordinate and $\varepsilon^2=0$. Reading the two sorts --- base $X$ and
tangent $v$ --- explicitly,
\begin{equation}\label{eq:Ttwosort}
   Tp(X,v) \;=\; p(X) + \varepsilon v\cdot\partial p(X)
            \;=\; \sum_a X^{E_a} + \varepsilon v\sum_a\sum_{e\in E_a} X^{E_a\setminus\{e\}} .
\end{equation}
We write $(-)\otimes W$, not ``$(-)\lhd D$'': substituting $D=y+\varepsilon v$ and truncating is
only the first-order shadow of the base change, and right $\lhd$-tensoring by $D$ as an
endofunctor of all of $\Poly$ does not give a tangent structure. On the first-order layer the two
descriptions agree, which is all we use.

\paragraph{Cartesian tangent categories and parallelizability.}
Recall \cite{CockettCruttwell2014} that a tangent category is \emph{Cartesian} when it has finite
products and the tangent functor preserves them, $T(A\times B)\cong TA\times TB$, so that the
trivial bundle $A\times m$ and the fibrewise addition make sense. A \emph{differential object} is
an object $m$ whose tangent bundle is trivial and linear. Lanfranchi's definition specializes to
this setting.

\begin{definition}[Lanfranchi \cite{Lanfranchi2026}, Def.~4.18]\label{def:par}
In a Cartesian tangent category, an object $M$ is \emph{parallelizable} if there is a linear
isomorphism of additive bundles over $M$,
\[
   \nu\colon TM \;\xrightarrow{\ \cong\ }\; M\times m \qquad\text{over } M,
\]
for a differential object $m$, compatible with the zero sections and the vertical lift.
\end{definition}

Theorem~4.20 of \cite{Lanfranchi2026} proves that, for a group object $G$ of a Cartesian tangent
category, being parallelizable is equivalent to the Lie-theoretic data (tangential/differential
representability of the Lie functor); and Corollary~4.22 records that in a \emph{generalized}
Cartesian differential category every object is parallelizable. Both facts are used below only to
diagnose the port; we quote and do not reprove them.

\section{The tangent fibre is the free module on the directions}\label{sec:fibre}

Everything rests on one computation.

\begin{lemma}\label{lem:fibre}
For $T=(-)\otimes W$ and $p=\sum_a y^{E_a}$, the fibre of the tangent bundle
$\pi\colon Tp\to p$ over a point of shape $a$ is the free module on the direction set $E_a$, of
rank $|E_a|$.
\end{lemma}

\begin{proof}
A point of $p$ over $X$ is a pair $(a, f\colon E_a\to X)$. By \eqref{eq:Ttwosort} a point of $Tp$
over it consists of the same base point together with the data recorded by the
$\varepsilon v$-summand
$\partial p(X)=\sum_a\sum_{e\in E_a}X^{E_a\setminus\{e\}}$: a choice of a direction $e\in E_a$ to
make infinitesimal, carrying a tangent value. The zero section is the point with all tangent
values $0$, and fibrewise addition adds tangent values slot by slot. Hence the fibre over
$(a,f)$ is the free module on $E_a$ --- one generator per direction --- of rank $|E_a|$.
\end{proof}

This is a statement about the \emph{fibres} of $Tp$; it needs no Cartesian apparatus and is the
robust kernel of the whole analysis. The verification script
\texttt{scratch/lie\_parallelizable\_verify.py} confirms, across ten categories, that the fibre
rank over a shape-$a$ point equals $|E_a|$.

\section{The Cartesian port is ill-posed}\label{sec:refutation}

Definition~\ref{def:par} requires a Cartesian tangent category. We show that no way of making
$\partial$ a tangent structure provides one in which parallelizability is both defined and
nontrivial. This proves Theorem~\ref{thm:I-intro}.

Throughout we use the species (counting) shadow in which a polynomial functor is tracked by its
counting polynomial in $X$, and the categorical product of polynomial functors is the product of
counting polynomials: a point of $p\times q$ over $X$ is a pair of points, so
$y^A\times y^B = y^{A\sqcup B}$ and the counting polynomials multiply, $X^{|A|}\cdot X^{|B|}$.

\subsection*{(a) $T$ does not preserve the categorical product}

Write $p'=\partial p$. By \eqref{eq:T}, $Tp = p + \varepsilon v\,p'$. For the categorical product,
\[
   T(p\times q) \;=\; pq + \varepsilon v\,(p'q + pq'),
\]
since $\partial$ is a derivation, while the product of the two tangent bundles is
\[
   Tp\times Tq \;=\; (p+\varepsilon v\,p')(q+\varepsilon v\,q')
   \;=\; pq + \varepsilon v\,(p'q+pq') + \varepsilon^2 v^2\, p'q'.
\]
These differ by the second-order term $\varepsilon^2 v^2\,p'q'$, which is nonzero: for $p=q=y$ it
is $\varepsilon^2 v^2$ (coefficient $1$), and for $p=q=y^2$ it is $4X^2\varepsilon^2v^2$
(coefficient $4X^2$); see \texttt{scratch/lie\_product\_compat.py}. The base change
$T(p\times q)=(p\times q)\otimes W$ tensors by a \emph{single} copy of $W$, in which
$\varepsilon^2=0$; the product bundle $Tp\times Tq$ carries two independent infinitesimal
directions --- morally it lives over $W\otimes W$, where the product of the two nilpotents does
not vanish --- and the categorical product supplies no map collapsing them. Hence $T$ does not
preserve $\times$, and $(\Poly,\times,T)$ is not a Cartesian tangent category. (The product rule
of \eqref{eq:T} holds only modulo $\varepsilon^2$, i.e.\ in the base-change setting itself, which
is monoidal and not Cartesian; see (c).)

\subsection*{(b) With the categorical product, only constants are parallelizable}

Suppose nonetheless that $Tp\cong p\times m$ over $p$, with $\times$ the categorical product. In
the counting shadow $p\times m$ is the product $p\cdot m$, so a necessary condition is that the
``quotient'' $m = Tp/p$ be itself a polynomial functor:
\[
   m \;=\; \frac{Tp}{p} \;=\; \frac{p+\varepsilon v\,p'}{p} \;=\; 1 + \varepsilon v\,\frac{p'}{p}.
\]
This is a polynomial functor iff $p \mid p'$ in $\Z[X]$. For a non-constant container
$\deg p' = \deg p - 1 < \deg p$, so $p\nmid p'$ unless $p'=0$, i.e.\ unless $p$ is constant. Hence
the \emph{only} parallelizable objects for the categorical product are the constant functors.
Even the terminal object $y$ fails: $Ty = X + \varepsilon v$, while $y\times m = X\cdot m$ can
never equal $X+\varepsilon v$. The port returns a degenerate answer with no dependence on the
category $p$ at all.

\subsection*{(c) The genuine tangent structure is monoidal, not Cartesian}

By \cite{MacBethUniqueness,LanfranchiLemay2025}, $\partial$ becomes a genuine (non-degenerate)
tangent bundle precisely as the Weil base change $(-)\otimes W$ --- equivalently, on the first
order, right-tensoring by the infinitesimal object in the \emph{monoidal} category $(\Poly,\lhd,y)$.
This is a representable/monoidal tangent setting. The composition product $\lhd$ is not Cartesian:
it has no diagonal $p\to p\lhd p$ natural in $p$ (the maps $p\to p\lhd p$ are extra structure ---
precisely the categories). Definition~\ref{def:par} lives only in Cartesian tangent categories, so
it does not apply to $(-)\otimes W$ as it stands.

\subsection*{(d) The Cartesian home is a GCDC --- vacuous and mistyped}

The honest Cartesian tangent category of containers is the arity (generalized) Cartesian
differential category $\SPoly$, and its ring completion $\PolyZ$: objects are arities $n$,
morphisms are polynomial maps, $T(n)=2n=n\times n$, and $\partial$ is the differential
\emph{combinator} on morphisms \cite{MacBethTangent,MacBethCartan}. A Cartesian differential
category is in particular a generalized Cartesian differential category, so by
Lanfranchi's Corollary~4.22 \cite{Lanfranchi2026} \emph{every} object is parallelizable: the
notion carries no information. Worse for the conjecture, here a container $p=\sum_a y^{E_a}$ is a
\emph{morphism} $1\to1$, not an object, so ``$p$ is parallelizable'' does not even typecheck.

\medskip
\noindent In every realization, parallelizability for $\partial$ is undefined (not Cartesian,
(a),(c)), degenerate (nothing non-constant, (b)), or vacuous and mistyped (GCDC, (d)). The literal
port of Theorem~4.20 to container comonoids is ill-posed. $\qed$

\section{The surviving invariant, and the groupoid refutation}\label{sec:survivor}

What \emph{is} well-posed and nontrivial is the fibre-rank condition of Lemma~\ref{lem:fibre},
intrinsic to the base-change bundle. This proves Theorem~\ref{thm:II-intro}.

\begin{definition}\label{def:treg}
A container $p=\sum_a y^{E_a}$ is \emph{tangent-regular} if the fibres of $Tp$ are all isomorphic
to a single module $m$.
\end{definition}

This is the honest residue of parallelizability available in the representable setting: not a
global trivialization (which (b) forbids), but the weaker demand that the bundle have a
\emph{constant fibre type}.

\begin{theorem}\label{thm:treg}
$p$ is tangent-regular iff the out-degree $a\mapsto|E_a|$ is constant on $\Ob$.
\end{theorem}

\begin{proof}
By Lemma~\ref{lem:fibre} the fibre over a shape-$a$ point is the free module of rank $|E_a|$. Two
free modules are isomorphic iff their ranks are equal, and a common $m$ exists for all shapes iff
all $|E_a|$ are equal.
\end{proof}

\begin{theorem}[blindness to composition]\label{thm:blind}
Tangent-regularity depends only on the underlying polynomial functor $U(p)$ --- equivalently, only
on the out-degree multiset $\{|E_a|\}$ --- hence factors through the forgetful functor
$U\colon\Cat\simeq\Comon(\Poly,\lhd)\to\Poly$. Consequently ``tangent-regular $\Leftrightarrow$
groupoid'' is false in both directions.
\end{theorem}

\begin{proof}
Tangent-regularity is defined from the fibres of $Tp$, and by Lemma~\ref{lem:fibre} those depend
only on the direction sets $E_a$, i.e.\ on $U(p)$, never on the comultiplication $\delta$. We
exhibit the two failures.
\begin{itemize}
\item \emph{(Tangent-regular $\not\Rightarrow$ groupoid.)} The idempotent monoid $\{1,e\}$ with
$e^2=e$ has one object and underlying functor $y^2$, identical to that of the group $\Z/2$. The two
have the \emph{same} tangent bundle; both are tangent-regular (one object, so out-degree is
trivially constant); but $\Z/2$ is a groupoid and $\{1,e\}$ is not. Tangent-regularity cannot tell
them apart.
\item \emph{(Groupoid $\not\Rightarrow$ tangent-regular.)} The disjoint groupoid
$\Z/2\sqcup\Z/3$ has objects of out-degrees $2$ and $3$, so $U=y^2+y^3$ is not tangent-regular by
Theorem~\ref{thm:treg}, although it is a groupoid. This direction is airtight from
Lemma~\ref{lem:fibre} alone.
\end{itemize}
\end{proof}

Why was the conjecture doomed? Because $\partial$ reads only the directions $E_a$ --- how many
morphisms leave each object --- and never the composition $\delta\colon p\to p\lhd p$. ``Groupoid'',
``connected'', ``invertible'' are properties of $\delta$, invisible to any invariant that factors
through $U$. No such invariant can be equivalent to ``groupoid''. The groupoid half of the
conjecture was unprovable in principle, not merely unproved.

\begin{corollary}[the positive shadow of Lanfranchi]\label{cor:shadow}
Every out-regular category --- one with $|E_a|$ constant --- is tangent-regular. In particular a
connected groupoid with object set $\Ob$ and vertex group $G$ (so $\Hom(a,b)\cong G$ for all
$a,b$) has $|E_a|=\sum_b|\Hom(a,b)|=|\Ob|\cdot|G|$, constant, hence is tangent-regular; so is any
one-object category (a monoid) and any codiscrete category.
\end{corollary}

This is the faithful remnant of ``Lie group $\Rightarrow$ parallelizable''. Homogeneity is exactly
the left-translation regularity --- composing with a connecting isomorphism $b\to a$ bijects
$\Hom(a,-)\cong\Hom(b,-)$ --- that a connected groupoid supplies. But by Theorem~\ref{thm:blind} it
is strictly weaker than being a groupoid.

\paragraph{Formal certification.}
The arithmetic core of Theorems~\ref{thm:treg}--\ref{thm:blind} is machine-checked in Lean~4
\cite{MacBethLean}. For out-degrees $d\colon\mathrm{Fin}\,k\to\N$, the file proves
$\mathrm{TangentRegular}\,d \leftrightarrow \forall i\,j,\ d_i=d_j$ (a common fibre rank exists iff
the out-degrees agree), and that the decidable test is \emph{permutation-invariant} --- it factors
through the multiset $\mathrm{List}/\mathrm{Perm}$ --- which is the formal content of
composition-blindness, since a comonoid's $\delta$ is not recoverable from the out-degree multiset.
Both counterexamples ($\{1,e\}$ versus $\Z/2$; $\Z/2\sqcup\Z/3$) are discharged as decidable facts.
The reduction that the representable fibre is the free module on $E_a$ (Lemma~\ref{lem:fibre}) and
the functor $U$ are cited, not formalized; \texttt{\#print axioms} reports
$\{\texttt{propext},\texttt{Quot.sound}\}$.

\section{The moral, and one open refinement}\label{sec:conclusion}

The tangent geometry that the container derivative supports is \emph{fibre-rank geometry over the
discrete set of objects}. It reads the out-degree sequence --- the shape of a category --- and is
blind to composition. ``Parallelizable $\Leftrightarrow$ Lie object'' does not survive the move
from the Cartesian affine world of \cite{Lanfranchi2026} to the representable/monoidal $\lhd$-world
of containers: the trivialization it demands collapses everything but constants
(\S\ref{sec:refutation}), and the honest residue --- constant fibre type --- cannot see the
composition that ``groupoid'' depends on (\S\ref{sec:survivor}). This is the same dividing line
that governs the fullness and uniqueness theorems of this program --- representability sees the
underlying functor and not the category structure laid over it --- appearing here in its
cautionary, boundary-drawing form. It belongs on a slide as a warning: before porting a
Cartesian-tangent theorem to containers, check that the ambient structure is $\times$ and not
$\lhd$, and that $T$ is Cartesian.

There is a genuine remaining question, sharper than the one refuted here. Theorem~\ref{thm:I-intro}
kills \emph{bare} Cartesian parallelizability, but it says nothing about a trivialization built
from the comonoid structure itself --- one that \emph{does} see composition.

\paragraph{Open problem.} For a $\lhd$-comonoid $p$ (a small category $\mathcal C$), does the
comultiplication $\delta$ induce a canonical ``left-translation'' trivialization of the vertical
part of the tangent bundle $Tp$ --- translate every tangent direction to a chosen basepoint by
composing along a connecting morphism? We conjecture that such a trivialization exists exactly when
the translations are invertible, so that this is where \emph{groupoids} (and connectedness)
genuinely enter --- the composition-sensitive layer that tangent-regularity provably cannot see.
Making this precise requires the differential-bundle/vertical structure of $Tp$ and the comonoid
maps worked out together; it is the correct home for the groupoid intuition, and we leave it to a
later study.

\paragraph{Acknowledgement.} This note is the cautionary companion to the derivative-uniqueness arc
and was prepared for the Strathclyde mini-course \emph{Tangent Categories and Containers}. Nothing
here is a correction to \cite{Lanfranchi2026}, whose theorem is about Cartesian tangent categories;
the port fails precisely because the container derivative's tangent structure is not one.

\begin{thebibliography}{9}

\bibitem{AAGM2003}
M.~Abbott, T.~Altenkirch, N.~Ghani, C.~McBride,
\emph{Derivatives of containers},
Typed Lambda Calculi and Applications (TLCA 2003), LNCS 2701, Springer, 2003, 16--30.

\bibitem{AhmanUustalu2016}
D.~Ahman, T.~Uustalu,
\emph{Directed containers as categories},
Proc.\ MSFP 2016, EPTCS 207, 2016, 89--98; arXiv:1604.01187.

\bibitem{BCS2009}
R.~F.~Blute, J.~R.~B.~Cockett, R.~A.~G.~Seely,
\emph{Cartesian differential categories},
Theory and Applications of Categories \textbf{22} (2009), no.~23, 622--672.

\bibitem{CockettCruttwell2014}
J.~R.~B.~Cockett, G.~S.~H.~Cruttwell,
\emph{Differential structure, tangent structure, and SDG},
Applied Categorical Structures \textbf{22} (2014), no.~2, 331--417.

\bibitem{GambinoKock2013}
N.~Gambino, J.~Kock,
\emph{Polynomial functors and polynomial monads},
Math.\ Proc.\ Cambridge Philos.\ Soc.\ \textbf{154} (2013), no.~1, 153--192; arXiv:0906.4931.

\bibitem{Lanfranchi2026}
M.~Lanfranchi,
\emph{The Lie group--Lie algebra correspondence in tangent categories},
preprint, arXiv:2609.03449 (2026).

\bibitem{LanfranchiLemay2025}
M.~Lanfranchi, J.-S.~P.~Lemay,
\emph{Representable tangent structures for affine schemes},
preprint, arXiv:2505.09080 (2025).

\bibitem{MacBethTangent}
MacBeth,
\emph{The container derivative is a tangent structure: a fibrational positioning}
(Kodamai internal note, 2026).

\bibitem{MacBethUniqueness}
MacBeth,
\emph{A uniqueness theorem for the container derivative: $\partial$ is the only nontrivial
representable tangent structure on finite-support polynomial functors}
(Kodamai internal note, 2026).

\bibitem{MacBethCartan}
MacBeth,
\emph{The container derivative has no fiberwise negation: a boundary for the Cartan calculus of
polynomial functors}
(Kodamai internal note, 2026).

\bibitem{MacBethLean}
MacBeth,
\emph{Lean~4 certification: tangent-regularity equals out-degree homogeneity, and its
composition-blindness}
(\texttt{2026-10-06-tangent-regular-outdegree.lean}; axioms
$\{\texttt{propext},\texttt{Quot.sound}\}$), 2026.

\end{thebibliography}

\end{document}
